\documentclass[aps,prl,twocolumn,showpacs,superscriptaddress]{revtex4-2}
\usepackage{amsmath,amssymb,booktabs,xcolor,hyperref,amsthm}
\newtheorem{theorem}{Theorem}

\begin{document}

\title{Appendix KB1: The Carbon Hopf Symmetry Index\\
$S_H = \min(n_{\rm filled},n_{\rm empty})/4$ and the Half-Fill
Uniqueness Theorem for the $n=2$ Hopf Multiplet}

\author{Michel Robert Cabri\'e}
\affiliation{Independent Researcher, Victoria, Australia}
\email{ZeroFreeParameters@gmail.com}
\date{20 March 2026}

\begin{abstract}
This appendix derives the Hopf Bond Symmetry Index $S_H$ for all 
$n=2$ elements, proves the uniqueness of $S_H=1$ at carbon ($Z=6$), 
derives the bond compression formula $L(n_\pi)/L_0 = 1-n_\pi r_{\rm tet}$ 
from $\pi$-bond Hopf mode geometry, and tabulates the half-fill 
series ($S_H=1$ elements) across all shells. Supporting Letter~82.
\end{abstract}

\maketitle

\section{The Hopf Bond Symmetry Index}

For $n=2$ elements with shell capacity $C=8$:
\begin{equation}
    S_H = \frac{\min(n_{\rm filled},\,8-n_{\rm filled})}{4}
\end{equation}

\begin{center}
\begin{tabular}{lcccr}
\toprule
Element & Filled & Empty & min & $S_H$ \\
\midrule
Li & 1 & 7 & 1 & 0.25 \\
Be & 2 & 6 & 2 & 0.50 \\
B  & 3 & 5 & 3 & 0.75 \\
\textbf{C}  & \textbf{4} & \textbf{4} & \textbf{4} & \textbf{1.00} $\star$ \\
N  & 5 & 3 & 3 & 0.75 \\
O  & 6 & 2 & 2 & 0.50 \\
F  & 7 & 1 & 1 & 0.25 \\
Ne & 8 & 0 & 0 & 0.00 \\
\bottomrule
\end{tabular}
\end{center}

\begin{theorem}[Carbon Uniqueness]
$S_H=1$ is achieved uniquely by carbon ($Z=6$) in the $n=2$ shell.
\end{theorem}
\begin{proof}
$S_H=1$ requires $\min(n_{\rm filled},\,8-n_{\rm filled})=4$, 
i.e.\ $n_{\rm filled}=4$. Only $Z=6$ has exactly 4 outer electrons 
in the $n=2$ shell. \qed
\end{proof}

\section{Bond Compression Derivation}

A $\sigma$ bond shares 2 Hopfions along the bond axis. Each 
$\pi$ bond shares 2 Hopfions via $r_{\rm tet}$ modes perpendicular 
to the axis, compressing the internuclear distance by $r_{\rm tet}$:
\begin{equation}
    \frac{L(n_\pi)}{L_0} = 1 - n_\pi \times r_{\rm tet}
\end{equation}

\begin{center}
\begin{tabular}{lclr}
\toprule
Bond & $n_\pi$ & BCT ratio & Error \\
\midrule
C--C & 0 & $1$ & --- \\
C=C  & 1 & $1-r_{\rm tet} = 0.8876$ & $+2.0\%$ \\
C$\equiv$C & 2 & $1-2r_{\rm tet} = 0.7753$ & $-0.5\%$ \\
\bottomrule
\end{tabular}
\end{center}

\section{The Half-Fill Series}

$S_H=1$ at the half-fill point of every shell:

\begin{center}
\begin{tabular}{rclll}
\toprule
$n$ & $Z$ & Symbol & Name & BCT significance \\
\midrule
2 & 6  & C  & Carbon    & Carbon-based life \\
3 & 14 & Si & Silicon   & Mineral structures \\
4 & 32 & Ge & Germanium & Semiconductor \\
5 & 50 & Sn & Tin       & BCT prediction \\
\bottomrule
\end{tabular}
\end{center}

\begin{thebibliography}{99}
\bibitem{L82} M.~R.~Cabri\'e, \textit{BCT Letter 82}, Zenodo (2026).
\bibitem{L78} M.~R.~Cabri\'e, \textit{BCT Letter 78}, Zenodo (2026).
\end{thebibliography}

\end{document}
